\section{Discussion}
\label{sec:discussion}

We set out to measure whether LLMs answer air pollution policy questions as
reliably for the Global South as for the Global North, whether prompt framing
changes the answer, and what mechanism drives any gap. Across 25 countries, 14
models, four languages, five tasks and two personas, scored against official
ground truth by code where the answer is a register value and by a three-vendor
panel elsewhere, the harms we can demonstrate are practical rather than
aetiological: the
three remedies an agency would reach for all fail, the most intuitive of them
makes the answer materially worse, and models are least reliable on the recall of
the binding value in force. A North/South tier difference persists, and we
identify the channel behind it only where the design has the resolution to do so.

\subsection{Three intuitive remedies, all failing --- and one that harms}

The local language does not help and materially hurts. Asking in the country's own
language lowers accuracy by $\nnativa$~pp, significantly in all three languages and
with the country as the unit of inference, and
the penalty scales with the scarcity of the language's corpus, with India --- the
highest-scoring country of all 25 in English --- falling furthest in Hindi. This
is the mechanism the resource hierarchy
predicts~\citep{joshi2020state,xue2021mt5}, and the measurement has the resolution
to see it because the unit is the language rather than the country. The effect
depends on no single country, model, task, persona condition or scoring
instrument, nor on translation quality, and it holds on an outcome involving no
judge at all: the model becomes markedly less likely to produce a checkable number
in the first place.

The local-manager persona does not narrow the gap (DiD $\ndid$~pp, permutation
$p=\npersonap$), consistent with evidence that expert personas do not reliably improve
factual accuracy~\citep{zheng2024helpful,mollick2025personas}. This null carries
more weight than a null about a folk habit would, because role framing is the
first thing every major vendor tells its users to do --- Anthropic, OpenAI, Google
and xAI converge on the same advice (Section~\ref{sec:related}), so an
administration following its supplier's documentation writes the sentence we
tested. What that documentation claims for the sentence is control over tone,
format and framing, and nothing in it marks the boundary: nothing tells a reader
that the technique which reliably shapes register does not also shape factual
reliability. Our estimate separates the two. A manipulation check confirms half
the responses explicitly take up the role, so the frame is received and changes
the voice; it does not change what the model knows. We tested it typed into the
message rather than installed as a system instruction, the channel a public
servant in a chat interface actually has (Section~\ref{sec:limits}).

The regional model is the \emph{worst} performer of all fourteen
($\delta=\nregdelta$), and it also trails its scale-matched peer on the
Portuguese-language Brazilian prompts it was tuned for ($\nhthreebra$ against $\nhthreebrallama$
for Llama~3.1~8B, on 20 cells each), so reaching for a locally tuned small model
trades away the scale and frontier training that carry domain accuracy without
buying back the local knowledge it promises. Two limits bound this: the narrow
test is small, and the model we could run is a community fine-tune of a 7B base,
not a commercial regional frontier model of the kind an agency might actually
procure --- for those the finding is a prediction, not a result, and the
released protocol is how to test it.

Practitioners deploy all three as fixes; none works, and the one a manager is most
likely to try first does active damage. For Global South environmental agencies,
no cheap prompt- or model-level workaround substitutes for verifying the binding
value against the official source.

\subsection{The recall floor sits on the values that bind}

The risk falls unevenly. A manager drafting prose about air quality policy is
comparatively well served; a manager retrieving the regulatory cutoff a decision
turns on is in the worst cell of the design (Section~\ref{sec:results:task}),
and worse still if they ask in a lower-resource language. The floor is not an
artefact of scoring: both recall tasks are decided by code against the register,
and a third of the wrong answers on the national standard cite a value from the
country's own official ladder rather than an invented one --- the models know the
ladder and not which rung binds.

\subsection{A tier difference we can show, and a gradient we cannot}

The geographic finding is a difference between groups, not a slope. The
North/South gap on the composite is modest, is the same size inside the
extension wave and on the pre-specified fifteen, and is concentrated: it scales
with how much the answer depends on a fact specific to the country, and
disappears on the one task with no register value to get right
(Section~\ref{sec:results:h1}). On the code-adjudicated outcomes alone, where no
judge can have contributed a deficit of its own, it is larger than on the
composite. The development gradient stays below the pre-specified threshold on
every reading, and that null is informative because the design has $\npower\%$ power
at $\rho=0.55$: a gradient of the size we committed to is excluded, a weak one is
not. The country ordering shows why the two results differ --- India tops all 25
countries despite being Global South, Canada and France sit low --- so the
disadvantage is a tier difference layered over country-specific idiosyncrasy
that a rank correlation over 25 countries cannot see through.

\subsection{The mechanism, where the design has resolution}

The pre-specified mechanism test fails, and the reason is a defect in the
instrument rather than an absence of signal. The proxy assigns the size of the
Wikipedia edition in a country's dominant official language, so nine of our 25
countries receive the value of the \emph{English} Wikipedia. The variable does not
measure how much text exists in a country's language; it measures whether the
country has English as an official language, which is a fact about colonial
history. We report this because it bears on any study that has used or will use
the same proxy.

Partial resolution comes from changing the unit of variation. Within the nine
anglophone countries the linguistic channel is not merely controlled but identical
by construction, and country coverage still predicts the specificity deficit
($\beta=\nhfouranglo$, $p=\nhfouranglop$): what that coefficient measures cannot be language.
What it can be is development. Entered together with coverage, HDI$\times$task
dependence takes the signal ($\beta=\nhfourjhdi$, $p=\nhfourjhdip$) and coverage
falls to $\beta=\nhfourjcob$ ($p=\nhfourjcobp$), because across these 25 countries the two
correlate at $\rho=\nrhocobhdi$ and the design has no lever that moves one without the
other. The mechanism is therefore identified only up to that pair: the deficit is
smaller where more about the country is on record, and whether the operative
record is the encyclopedic footprint or the wider apparatus of a developed state
this study cannot say.

\subsection{The methodological lesson: the key matters more than the sample}

The same scored cells support different substantive conclusions depending on
the instrument used to score them, and the size of that difference is the
lesson. Scoring two tasks against placeholder keys and delegating range
arithmetic to a judge yielded a significant development gradient and a small
language penalty; rebuilding the keys from official registers and moving the
comparison into code halved the gradient and more than doubled the penalty
(Section~\ref{sec:results:h1}, \ref{sec:results:h2}). Moving the third
register-valued task from two collection judges to code changed its headline
odds ratio by a third, because the judges had been split along the same line as
the tiers. And the key itself kept moving under scrutiny: the United Kingdom
entry carried a 2040 target in place of the limit in force, and the Bangladesh
entry carried the 2005 value after the 2022 Rules had raised it, so that a model
citing the correct current standard for a Global South country was being scored
wrong until an external reader caught it. The dataset never changed; only the
instrument did. The field's instinct when a geographic-bias effect proves
fragile is to add countries, but the ground truth deserves the scrutiny first: a
defective key does not merely add noise, it adds \emph{structured} noise,
because the defects concentrate wherever official texts are hardest to obtain or
most recently revised --- which is exactly where the hypothesis lives. The
released registry therefore carries a content hash (Supplementary Material) so
that any reader can tell which version of the key a number was computed against.

\subsection{Why the domain matters}

Fine particulate matter is among the leading environmental risk factors for
premature mortality, and the burden falls disproportionately on the Global South:
the World Health Organization attributes \num{88548} deaths in Brazil in 2021 to
ambient air pollution (uncertainty interval
\numrange{69477}{108431})~\citep{who2026gho_air41}, and short-term exposure
continues to drive measurable mortality in Brazilian
cities~\citep{requia2024shortterm}. When a manager delegates a normative lookup or
an episode recommendation to an LLM, an inaccurate answer can propagate into a
regulatory brief or an emergency response.

We do not claim the model is least reliable where consequences are greatest,
because we did not measure it: disease burden is not among our country covariates,
and the one country-level observation we have points the other way, since India
carries one of the largest PM\textsubscript{2.5} mortality burdens in the world
and is also the highest-scoring country in our sample. Whether reliability is
misaligned with need is answerable by adding an age-standardised
attributable-mortality covariate, but not by this design. We interpret the pattern
we \emph{did} measure through the coloniality-of-knowledge
framework~\citep{mohamed2020decolonial}: LLMs are compressed summaries of corpora
produced under uneven digital and publication access, and the language penalty is
a measurable link from corpus asymmetry to knowledge asymmetry.

\subsection{Limitations and future opportunities}
\label{sec:limits}

Each limit below bounds this design, and each names the study that would remove
it. We list them in the order we would run them.

\emph{Role framing was tested in the user turn, not the system prompt.} The
manager frame is prefixed to the message rather than installed as a system
instruction, where three of the four vendor guides place it --- Google's admits
either that slot or the start of the user prompt, which is the one we used. The
choice follows the
population we care about --- public servants in the chat interface, who have no
system prompt to configure --- but our null bounds that user's channel, not a
developer's. Whether a system-instruction persona recovers factual accuracy where
a typed one does not is a clean two-cell experiment on the same items, and it is
the first study we would run.

\emph{From panel-scored to gold-anchored.} No human expert re-scored responses
here. Three properties bound what that costs: where the answer is a register value
the verdict is computed by code; where judgement is required the score is the mean
of three vendors rather than one; and every reported effect is a between-group
contrast, in which a judge's systematic severity cancels. Restricting to objective
factual accuracy \emph{strengthens} both headline effects, the opposite of what a
permissive judge would produce. A blinded human-expert validation of a stratified
sample --- roughly 150 responses scored by three raters, which the released
pipeline exports directly --- would convert an official-source-anchored result
into a gold-standard-anchored one.

\emph{From a within-country signal to a causal design.} Within countries we remove
country-level confounding without removing task-by-country confounding. A design
that varies country coverage while holding development fixed, or a longitudinal
study tracking accuracy as a country's encyclopedic footprint grows, would close
that gap.

\emph{From a null gradient to an adequately powered test.} The tier difference is
established while the gradient is not, at $n=25$. A pre-registered replication at
substantially larger $n$, which the open dataset makes cheap, would settle whether
the gradient is genuinely absent or still underpowered.

\emph{From three native languages to the full multilingual matrix.} The strongest
result rests on three native languages, so its \emph{mechanism} is reported
descriptively even though the effect is not. The translation-quality objection is
settled empirically by the back-translation census reported above; what remains
open is breadth, not validity, and scaling to more languages is the highest-value
extension this study points to.

\emph{From one domain to a cross-domain program.} All effects come from one
high-stakes applied domain and one corpus family (Wikimedia), and we confine the
mechanistic language accordingly. The same task suite and ground-truth-registry
method transfer directly to other regulated, source-anchored domains --- water
quality, food safety, vaccine policy.

\emph{Served-model stack as the unit of analysis.} We evaluate models \emph{as
served} through fixed access stacks, which is what a public manager queries; a
residual gap could in principle reflect the stack rather than the weights. A
controlled quasi-isolation design with identical open weights across stacks would
separate the two, which the per-call provenance we release already enables.

\subsection{What this study contributes}

The three fixes practitioners reach for first fail together:
the local language actively harms, the persona frame does nothing, and the
regionally fine-tuned 7B model is the single worst performer of fourteen. A benchmark's most
useful output is often telling the field what does not work, with effect sizes
behind it. The methodological result
is that in this literature the ground-truth key, not the sample size, is the
binding constraint. And the design separates what 25 countries can demonstrate
from what they cannot, reporting the second set as open rather than as supported.